\clearpage
\section{Miền tin cậy và khoảng tin cậy đồng thời}
\label{sec:confidence}

\subsection{Ma trận nào mô tả quan sát, ma trận nào mô tả trung bình?}
\label{sec:covariance_roles}
Trước khi dựng miền tin cậy, cần biết ma trận đang đo độ phân tán của
\emph{đại lượng nào}. Các ma trận dưới đây có cùng kích thước $p\times p$
nhưng không có cùng vai trò.
\begin{center}
\small
\begin{tabular}{L{.19\linewidth}L{.25\linewidth}L{.44\linewidth}}
\toprule
Ma trận & Đại lượng được mô tả & Ý nghĩa và cách sử dụng\\
\midrule
$\Sigma$ & Một quan sát $\mathbf X_j$ & Hiệp phương sai tổng thể;
 mô tả đám mây quan sát, thường chưa biết.\\[3pt]
$S$ & Ước lượng $\Sigma$ & Hiệp phương sai mẫu không chệch;
 phản ánh độ phân tán và liên hệ giữa các cột dữ liệu.\\[3pt]
$\Sigma/n$ & Trung bình mẫu $\Xbar$ & Hiệp phương sai thực của trung bình
 mẫu khi các quan sát iid.\\[3pt]
$S/n$ & Ước lượng $\Cov(\Xbar)$ & Mô tả sai số của trung bình;
 nghịch đảo của nó tạo thước đo trong $T^2$.\\[3pt]
$\widehat\Sigma$ & MLE của $\Sigma$ & Bằng $(n-1)S/n$, phục vụ hàm hợp
 lý; không phải hiệp phương sai của trung bình mẫu.\\
\bottomrule
\end{tabular}
\end{center}
\nguon{Chương 4, §4.3, Result 4.11, tr.~171--172; §4.4, (4-23), tr.~174}

\paragraph{Đường chéo: biến thiên của từng biến riêng.}
$s_{ii}$ ước lượng phương sai của biến thứ $i$; $s_{ii}/n$ ước lượng
phương sai của trung bình thành phần. Do đó sai số chuẩn là
$\sqrt{s_{ii}/n}$. Không dùng $\sqrt{s_{ii}}$ làm sai số chuẩn của
trung bình: nó mô tả độ phân tán của quan sát.

\paragraph{Ngoài đường chéo: cách hai biến biến thiên cùng nhau.}
$s_{ik}$ mang thông tin ghép cặp giữa hai biến. Tương quan mẫu là
$r_{ik}=s_{ik}/\sqrt{s_{ii}s_{kk}}$. Nếu chia toàn bộ ma trận cho $n$,
tương quan không đổi vì hệ số $1/n$ triệt tiêu ở tử và mẫu.
Lấy nhiều quan sát giúp ước lượng trung bình chính xác hơn, nhưng không
đồng nghĩa hai thành phần trung bình mẫu trở nên không tương quan.

\paragraph{Ma trận nghịch đảo: độ lệch theo hướng nào đáng chú ý?}
Hotelling dùng $T^2=\mathbf r'(S/n)^{-1}\mathbf r$ với
$\mathbf r=\xbar-\vmu_0$. Phép nghịch đảo hiệu chỉnh đồng thời đơn vị,
phương sai và hiệp phương sai. Sau khi xoay theo các vectơ riêng,
\[
 T^2=n\sum_{i=1}^p\frac{(\mathbf e_i'\mathbf r)^2}{\ell_i},
 \qquad S\mathbf e_i=\ell_i\mathbf e_i.
\]
Độ lệch theo hướng có phương sai nhỏ $\ell_i$ bị đánh trọng số lớn hơn.
Đó là lý do một điểm gần tâm theo khoảng cách thường vẫn có thể nằm
ngoài elip. Nếu một trị riêng gần 0, nghịch đảo rất nhạy với sai số dữ
liệu; nếu trị riêng bằng 0, công thức cổ điển không dùng được.

\hieu{Hiệp phương sai không chỉ quyết định khoảng riêng của từng biến.
Nó còn cho biết các \emph{cặp giá trị trung bình} nào phù hợp với nhau.
Miền elip giữ thông tin về các hướng đó; một hình chữ nhật dựng từ hai
khoảng thành phần chỉ đặt giới hạn riêng cho từng tọa độ.}


\subsection{Đảo kiểm định để tìm các giá trị trung bình còn phù hợp}
Ở phần trước, ta kiểm tra một vectơ $\vmu_0$ cụ thể. Bây giờ thử mọi vectơ
$\vmu$ và giữ các vectơ không bị bác bỏ. Ta thu được miền
\begin{equation}
 \mathcal C=\left\{\vmu:
 n(\xbar-\vmu)'S^{-1}(\xbar-\vmu)\le c_T^2\right\},
 \qquad c_T^2=\frac{p(n-1)}{n-p}F_{p,n-p;1-\alpha}.
 \label{eq:ellipsoid}
\end{equation}
Dưới giả thiết mẫu iid chuẩn và $n>p$, đây là miền tin cậy mức
$1-\alpha$ cho vectơ trung bình.
\nguon{Chương 5, §5.4, miền tin cậy (5-18), tr.~221; diễn giải hình học tr.~221--223}

\hieu{Tâm của miền là trung bình mẫu. Vectơ trung bình nào quá xa tâm
so với sai số lấy mẫu bị loại. Miền thể hiện các giá trị \emph{đồng thời}
của tất cả các thành phần, có xét đến tương quan.}

Nếu $\ell_i$ là trị riêng của $S$ với vectơ riêng đơn vị $\mathbf e_i$,
bán trục theo hướng $\mathbf e_i$ dài $\sqrt{c_T^2\ell_i/n}$.
Phân biệt với elipsoid xác suất cho quan sát trong
Phần~\ref{sec:normal}: ở đây $S$ ước lượng từ mẫu và có hệ số $1/n$.

Nếu $\vmu_0\notin\mathcal C$, kiểm định mức $\alpha$ bác bỏ $H_0$;
nếu $\vmu_0\in\mathcal C$, kiểm định chưa bác bỏ. Đây là cùng một quy tắc,
không phải hai kết quả suy diễn độc lập cần chỉnh mức ý nghĩa lần nữa.

\subsection{Khoảng tin cậy cho một tổ hợp bất kỳ}
Cho $\mathbf a$ là vectơ hệ số. Từ miền \eqref{eq:ellipsoid}, các giá trị
có thể của $\mathbf a'\vmu$ nằm trong
\begin{equation}
 \mathbf a'\xbar\ \pm\
 \sqrt{c_T^2}\sqrt{\frac{\mathbf a'S\mathbf a}{n}}.
 \label{eq:t2interval}
\end{equation}
\begin{result}
Với mẫu iid chuẩn, $\Sigma$ xác định dương và $n>p$, các khoảng
\eqref{eq:t2interval} chứa $\mathbf a'\vmu$ \emph{đồng thời cho mọi}
$\mathbf a$, với xác suất $1-\alpha$ trước khi lấy mẫu.
\end{result}
\nguon{Chương 5, §5.4, Result 5.3, tr.~225}

Trong công thức, $\sqrt{\mathbf a'S\mathbf a/n}$ là sai số chuẩn của
$\mathbf a'\xbar$, còn $\sqrt{c_T^2}$ là hệ số mở rộng để bảo đảm đồng
thời cho \emph{mọi hướng}. Nó thường lớn hơn hệ số của một khoảng $t$
cho một tổ hợp cố định.

Vì đã bảo vệ mọi hướng, kết quả còn cho phép xem dữ liệu rồi chọn hướng
đáng chú ý để khảo sát. Điều này khác Bonferroni thông thường cho một
danh sách hữu hạn đã định trước. Phụ lục~\ref{app:projection} giải thích
bằng bất đẳng thức Cauchy--Schwarz.

\subsection{Trung bình thành phần và hiệu hai trung bình}
Chọn $\mathbf a=\mathbf e_i$ (vectơ có 1 ở vị trí $i$ và 0 ở nơi khác):
\begin{equation}
 \mu_i:\quad\bar x_i\pm\sqrt{c_T^2}\sqrt{s_{ii}/n}.
 \label{eq:t2components}
\end{equation}
Chọn $\mathbf a=\mathbf e_i-\mathbf e_k$:
\begin{equation}
 \mu_i-\mu_k:\quad
 (\bar x_i-\bar x_k)\pm\sqrt{c_T^2}
 \sqrt{\frac{s_{ii}+s_{kk}-2s_{ik}}n}.
 \label{eq:t2difference}
\end{equation}
\nguon{Chương 5, §5.4, (5-24), tr.~225; (5-25), tr.~226}

\luuy{Đây là hiệu hai \emph{thành phần đo trên cùng các đối tượng},
không phải hiệu trung bình của hai mẫu độc lập. Bỏ $-2s_{ik}$ là bỏ
thông tin ghép cặp. Đồng thời, chỉ diễn giải một hiệu nếu đơn vị đo và
câu hỏi thực tế làm hiệu đó có ý nghĩa.}

\begin{example}[Một bộ số minh họa cho các khoảng]
Giả sử mẫu iid chuẩn hai chiều có $n=40$, $\xbar=(10,20)'$ và
\begin{equation}
 S=\begin{pmatrix}4&3\\3&9\end{pmatrix},\qquad\alpha=0.05.
 \label{eq:intervaldata}
\end{equation}
Đây là \emph{số tóm tắt giả định để minh họa}, không phải dữ liệu của mẫu
$n=4$ ở phần trước. Ta có
\[
 c_T^2=\frac{2\cdot39}{38}F_{2,38;0.95}\approx6.6604,
 \qquad\sqrt{c_T^2}\approx2.5808.
\]
Ba sai số chuẩn là
\[
 \sqrt{4/40}=0.3162,\quad\sqrt{9/40}=0.4743,\quad
 \sqrt{(4+9-6)/40}=\sqrt{7/40}=0.4183.
\]
Các khoảng $T^2$ đồng thời:
\begin{align*}
 \mu_1 &: 10\pm0.8161=[9.1839;10.8161],\\
 \mu_2 &: 20\pm1.2242=[18.7758;21.2242],\\
 \mu_1-\mu_2 &: -10\pm1.0796=[-11.0796;-8.9204].
\end{align*}
Khoảng hiệu không chứa 0, nên có bằng chứng hai trung bình thành phần
khác nhau, khi hiệu có ý nghĩa về đơn vị và thiết kế đo.
\end{example}

\subsection{Hình chiếu của miền elip và ý nghĩa của hình chữ nhật}
Các khoảng cho từng thành phần là \emph{hình chiếu} của elip lên các trục.
Hình chữ nhật tạo bởi chúng chứa toàn bộ elip và thêm các góc ngoài elip.
Do đó việc từng thành phần của $\vmu_0$ nằm trong khoảng riêng $T^2$
\emph{chưa đủ} để nói vectơ $\vmu_0$ nằm trong elip.

\begin{figure}[htbp]
\centering
\begin{tikzpicture}[x=3.4cm,y=2.3cm,>=Stealth]
 \draw[->] (-1.18,0)--(1.2,0) node[right]{$\mu_1-10$};
 \draw[->] (0,-1.6)--(0,1.6) node[above]{$\mu_2-20$};
 \draw[dashed] (-.816114,-1.224171) rectangle (.816114,1.224171);
 \draw[domain=0:360,samples=120,smooth,variable=\t]
 plot ({.816114*cos(\t)},{.612085*cos(\t)+1.060163*sin(\t)});
 \fill (0,0) circle (.025) node[below left]{tâm};
 \fill (.75,-1.1) circle (.025) node[right]{$Q$};
 \node[below] at (.816114,0) {$0.8161$};
 \node[above left] at (0,1.224171) {$1.2242$};
\end{tikzpicture}
\caption{Elip 95\% của ví dụ $n=40$, vẽ theo độ lệch khỏi trung bình mẫu.
Nét đứt là hình chữ nhật từ các khoảng thành phần $T^2$.
Điểm $Q$ nằm trong hình chữ nhật nhưng ngoài elip.}
\label{fig:projection}
\end{figure}

Trong hình, $Q$ biểu diễn $\vmu_0=(10.75,18.9)'$.
Hai độ lệch nằm trong các nửa độ rộng $0.8161,1.2242$,
nhưng
\[
 40(0.75,-1.1)S^{-1}(0.75,-1.1)'\approx22.0037>6.6604.
\]
Kiểm định vectơ trung bình vẫn bác bỏ giả thuyết này. Điểm đó lệch ngược
chiều tương quan dương, một hướng ít biến thiên của dữ liệu.

\subsection{Bonferroni: bảo vệ một danh sách hữu hạn}
Giả sử trước khi xem dữ liệu, ta muốn ước lượng $m$ đại lượng
$\theta_j=\mathbf a_j'\vmu$, $j=1,\ldots,m$.
Một khoảng $t$ cho mỗi đại lượng ở mức $1-\alpha/m$ là
\begin{equation}
 \mathbf a_j'\xbar\ \pm\
 t_{n-1;1-\alpha/(2m)}\sqrt{\frac{\mathbf a_j'S\mathbf a_j}{n}}.
 \label{eq:bonf}
\end{equation}
Dưới mẫu iid chuẩn, mỗi khoảng có xác suất không phủ bằng $\alpha/m$.
Bất đẳng thức hợp cho
\[
 \Prob(\text{ít nhất một khoảng không phủ})
 \le\sum_{j=1}^m\frac\alpha m=\alpha.
\]
Vậy xác suất tất cả các khoảng đồng thời phủ các tham số tương ứng
\emph{ít nhất} $1-\alpha$. Không cần các khoảng độc lập.
\nguon{Chương 5, §5.4, bất đẳng thức (5-28) và khoảng (5-29), tr.~232;
 mở rộng sang tổ hợp và so độ dài (5-30), tr.~234}

\hieu{Ta chia ngân sách sai số tổng $\alpha$ cho $m$ phát biểu.
Mỗi phát biểu được phép sai ít hơn, nên khoảng riêng phải rộng hơn khoảng
95\% thông thường. Đổi lại, cả nhóm có bảo đảm đồng thời.}

Với hai trung bình và một hiệu, có $m=3$, dù dữ liệu chỉ có $p=2$ biến.
Nếu chỉ báo cáo hai trung bình, $m=2$. Không mặc định $m=p$ nếu danh sách
có thêm các tổ hợp khác. Hệ số dùng $\alpha/(2m)$ vì mỗi khoảng là hai phía.

\begin{example}[So sánh trên cùng danh sách ba phát biểu]
Với \eqref{eq:intervaldata}, danh sách là
$\mu_1,\mu_2,\mu_1-\mu_2$, nên $m=3$.
$t_{39;1-0.05/6}\approx2.5017$.
\begin{center}
\begin{tabular}{lrr}
\toprule
Đại lượng & Nửa độ rộng $T^2$ & Nửa độ rộng Bonferroni\\
\midrule
$\mu_1$ & 0.8161 & 0.7911\\
$\mu_2$ & 1.2242 & 1.1866\\
$\mu_1-\mu_2$ & 1.0796 & 1.0465\\
\bottomrule
\end{tabular}
\end{center}
Các khoảng Bonferroni là $[9.2089;10.7911]$,
$[18.8134;21.1866]$ và $[-11.0465;-8.9535]$.
Trong ví dụ này Bonferroni hẹp hơn, nhưng chỉ bảo vệ danh sách đã xác định,
còn $T^2$ bảo vệ mọi tổ hợp tuyến tính.
\end{example}

\subsection{Dựng và nhận xét miền elip từng bước}
\label{sec:construct_regions}
Trong hai chiều, đặt $\delta_i=\mu_i-\bar x_i$ và
$r=s_{12}/\sqrt{s_{11}s_{22}}$, với $S\succ0$ nên $|r|<1$.
Ta có
\[
 S^{-1}=\frac1{s_{11}s_{22}-s_{12}^2}
 \begin{pmatrix}s_{22}&-s_{12}\\-s_{12}&s_{11}\end{pmatrix}.
\]
Phương trình miền elip viết tường minh thành
\begin{equation}
 \frac n{1-r^2}\left[
 \frac{\delta_1^2}{s_{11}}
 -\frac{2r\delta_1\delta_2}{\sqrt{s_{11}s_{22}}}
 +\frac{\delta_2^2}{s_{22}}\right]\le c_T^2.
 \label{eq:ellipse_correlation}
\end{equation}
Hạng $\delta_1\delta_2$ cho thấy miền phụ thuộc cách \emph{hai tọa độ
cùng lệch}. Nó không thể nói đầy đủ bằng hai bất đẳng thức riêng lẻ.

\buoc{1}{Xác định tâm và mức tin cậy}
Tâm là $\xbar$, không phải $\vmu_0$. Chọn $1-\alpha$, tính $c_T^2$
theo $n,p$. Một giả thuyết cụ thể chỉ là một điểm để so với miền đã dựng.

\buoc{2}{Xác định các hướng và bán trục}
Giải $S\mathbf e_i=\ell_i\mathbf e_i$, chuẩn hóa $\mathbf e_i$ về độ
dài 1. Bán trục theo hướng đó là $a_i=\sqrt{c_T^2\ell_i/n}$.
Hai đầu mỗi trục là $\xbar\pm a_i\mathbf e_i$.
Trong hai chiều có thể tính trị riêng trực tiếp:
\[
 \ell_{\pm}=\frac{s_{11}+s_{22}
 \pm\sqrt{(s_{11}-s_{22})^2+4s_{12}^2}}2.
\]
Trị riêng lớn cho trục dài: hướng ấy có độ bất định lớn hơn.
Khi $s_{12}=0$, trục elip song song trục tọa độ. Nếu phương sai khác
nhau, elip vẫn không phải đường tròn.
\nguon{Chương 5, §5.4, (5-18), (5-19), tr.~221; hình học dựa trên (4-7), tr.~153}

\begin{example}[Trục của elip trong ví dụ $n=40$]
Với \eqref{eq:intervaldata}, $S=\begin{pmatrix}4&3\\3&9\end{pmatrix}$,
hai trị riêng là $(13\pm\sqrt{61})/2$, xấp xỉ 10.4051 và 2.5949.
Các bán trục 95\% dài khoảng 1.3163 và 0.6573.
Hướng trục dài có thể chọn $\mathbf e_+=(0.4242,0.9056)'$,
nghiêng khoảng $64.9^\circ$ so với chiều dương trục $\mu_1$.
Tương quan dương làm hướng dài nghiêng lên; độ nghiêng không nhất thiết
$45^\circ$ vì hai phương sai 4 và 9 khác nhau.
\end{example}

\buoc{3}{Kiểm tra một điểm và viết nhận xét}
Với điểm ứng viên $\vmu_0$, tính
$q=n(\xbar-\vmu_0)'S^{-1}(\xbar-\vmu_0)$.
Nếu $q\le c_T^2$, điểm nằm trong hoặc trên biên; nếu $q>c_T^2$, điểm
nằm ngoài. Nhận xét nên nêu tâm, hướng biến thiên, độ dài trục và vị trí
của điểm, rồi nối với quyết định Hotelling ở cùng mức $\alpha$.

Một câu nhận xét đủ ý là: ``Elip tập trung quanh trung bình mẫu và kéo
dài theo hướng hai biến cùng tăng. Điểm giả thuyết nằm ngoài elip vì lệch
theo hướng ít biến thiên; do đó kiểm định Hotelling bác bỏ ở mức 5\%.''
Không viết ``điểm nằm ngoài nên dữ liệu chắc chắn sai'', và không gán
xác suất 95\% cho một tham số cố định nằm trong miền đã quan sát.

\subsection{Dựng miền chữ nhật và phân biệt hai cách dựng}
\label{sec:rectangles}
Với hai trung bình thành phần, hình chữ nhật từ hình chiếu $T^2$ là
\begin{equation}
 \mathcal R_T=\prod_{i=1}^2
 \left[\bar x_i-\sqrt{c_T^2}\sqrt{s_{ii}/n};
 \bar x_i+\sqrt{c_T^2}\sqrt{s_{ii}/n}\right].
 \label{eq:rectangle_T}
\end{equation}
Nó chứa toàn bộ elip $\mathcal C$. Bảo đảm phủ của hình chữ nhật thành
phần là \emph{ít nhất} $1-\alpha$; kết quả phủ đúng $1-\alpha$ cho
\emph{mọi tổ hợp} không đồng nghĩa chỉ hai khoảng thành phần cũng có
xác suất phủ đúng bằng mức đó.

Nếu chỉ hai trung bình thành phần được định trước, Bonferroni dùng $m=2$:
\begin{equation}
 \mathcal R_B=\prod_{i=1}^2
 \left[\bar x_i-k_B\sqrt{s_{ii}/n};
 \bar x_i+k_B\sqrt{s_{ii}/n}\right],\quad
 k_B=t_{n-1;1-\alpha/4}.
 \label{eq:rectangle_B}
\end{equation}
Tâm vẫn là $\xbar$; hai nửa cạnh lần lượt là
$h_{B1}=k_B\sqrt{s_{11}/n}$ và $h_{B2}=k_B\sqrt{s_{22}/n}$.
Các cạnh song song với trục tọa độ. Để kiểm tra điểm, chỉ cần
$|\delta_1|\le h_{B1}$ \emph{và} $|\delta_2|\le h_{B2}$.
\nguon{Chương 5, §5.4, hình chiếu của khoảng $T^2$, tr.~226;
 Bonferroni (5-28), (5-29), tr.~232; so sánh các miền trong Example 5.6, tr.~233--234}

\begin{example}[Hai miền không luôn chứa nhau]
Vẫn dùng $n=40$, $\xbar=(10,20)'$, $S$ trong \eqref{eq:intervaldata}.
Lần này danh sách Bonferroni \emph{chỉ có hai trung bình}, khác danh sách
$m=3$ ở ví dụ trước. Với $\alpha=0.05$,
\[
 k_B=t_{39;0.9875}\approx2.3313,\qquad
 (h_{B1},h_{B2})\approx(0.7372,1.1058).
\]
Vậy
\[
 \mathcal R_B=[9.2628;10.7372]\times[18.8942;21.1058].
\]
Xét hai điểm:
\begin{itemize}
 \item $P=(10.70,18.95)'$ có độ lệch $(0.70,-1.05)'$,
 nằm trong Bonferroni nhưng $q=19.6>6.6604$, nên nằm ngoài elip.
 \item $R=(10.80,20.60)'$ có độ lệch $(0.80,0.60)'$,
 có $q=6.4<6.6604$, nên nằm trong elip, nhưng $0.80>0.7372$,
 do đó nằm ngoài Bonferroni.
\end{itemize}
Như vậy cả hai miền đều có bảo đảm tin cậy thích hợp, nhưng hình dạng
và cách đánh giá một điểm khác nhau. Không khẳng định hình chữ nhật
Bonferroni luôn nằm trong elip hoặc luôn chứa elip.
\end{example}

\begin{figure}[htbp]
\centering
\begin{tikzpicture}[x=3.4cm,y=2.3cm,>=Stealth]
 \draw[->] (-1.18,0)--(1.2,0) node[right]{$\delta_1$};
 \draw[->] (0,-1.6)--(0,1.6) node[above]{$\delta_2$};
 \draw[dashed] (-.816114,-1.224171) rectangle (.816114,1.224171);
 \draw[densely dotted,thick] (-.737210,-1.105816) rectangle (.737210,1.105816);
 \draw[domain=0:360,samples=120,smooth,variable=\t]
 plot ({.816114*cos(\t)},{.612085*cos(\t)+1.060163*sin(\t)});
 \fill (.70,-1.05) circle (.025) node[above left]{$P$};
 \fill (.80,.60) circle (.025) node[right]{$R$};
 \fill (0,0) circle (.02) node[below left]{$\xbar$};
\end{tikzpicture}
\caption{Cùng dữ liệu và mức 95\%: nét liền là elip Hotelling,
 nét đứt là hình chữ nhật từ các khoảng thành phần $T^2$, nét chấm là
 Bonferroni cho $m=2$. Trục thể hiện độ lệch khỏi $\xbar$.}
\label{fig:three_regions}
\end{figure}

\begin{center}
\small
\begin{tabular}{L{.21\linewidth}L{.22\linewidth}L{.23\linewidth}L{.23\linewidth}}
\toprule
Tiêu chí & Elip Hotelling & Chữ nhật $T^2$ & Chữ nhật Bonferroni\\
\midrule
Tâm & $\xbar$ & $\xbar$ & $\xbar$\\[3pt]
Thông tin dùng để dựng & Toàn bộ $S$, kể cả $s_{12}$ & $s_{11},s_{22}$ và hệ số $T^2$ & $s_{11},s_{22}$ và hệ số theo $m$\\[3pt]
Hướng biên & Theo vectơ riêng & Song song trục tọa độ & Song song trục tọa độ\\[3pt]
Phủ trong mẫu chuẩn & Đúng $1-\alpha$ cho vectơ & Ít nhất $1-\alpha$ cho các thành phần & Ít nhất $1-\alpha$ cho danh sách định trước\\[3pt]
Nhận xét điểm & Dùng dạng toàn phương & Kiểm tra từng tọa độ & Kiểm tra từng tọa độ\\[3pt]
Quan hệ hình học & Luôn nằm trong $\mathcal R_T$ & Chứa elip & Không có quan hệ chứa nhau cố định với elip\\
\bottomrule
\end{tabular}
\end{center}
Điểm ngoài elip dẫn tới bác bỏ Hotelling. Điểm ngoài Bonferroni vi phạm
ít nhất một khoảng thành phần đã hiệu chỉnh; đó không phải cùng một
quy tắc kiểm định Hotelling, nên các kết luận có thể khác nhau.

\subsection{Vì sao chữ nhật Bonferroni không trực tiếp phản ánh tương quan?}
\label{sec:bonf_correlation}
Công thức \eqref{eq:rectangle_B} không chứa $s_{12}$: mỗi cạnh được dựng
từ một khoảng một chiều cho $\mu_i$. Sau đó Bonferroni kết hợp các
phát biểu bằng một bất đẳng thức đúng với \emph{mọi kiểu phụ thuộc}.
Phương pháp bảo đảm sai số mà không cần ước lượng xác suất chung chính
xác theo tương quan. Vì vậy hình chữ nhật không xoay theo đám mây dữ liệu.

Để thấy rõ, giữ $n=40$, $\xbar=(10,20)'$, mức 95\%, xét
\begin{equation}
 S_\rho=\begin{pmatrix}4&6\rho\\6\rho&9\end{pmatrix},\quad
 \rho=-0.5,\ 0,\ 0.5.
 \label{eq:covariance_family}
\end{equation}
Đây là ba ma trận hiệp phương sai mẫu giả định, cùng phương sai thành
phần nhưng tương quan khác nhau. Trong cả ba, hình chữ nhật Bonferroni
có cùng tâm, cùng nửa cạnh $(0.7372,1.1058)$.
Elip thì đổi hướng và bán trục; tại $\rho=0$, hai bán trục là 0.8161
và 1.2242, còn tại $\rho=\pm0.5$ là 0.6573 và 1.3163.

\begin{figure}[htbp]
\centering
\begin{tikzpicture}[>=Stealth,font=\small]
 \foreach \corr/\shift in {-0.5/0,0/4.6,0.5/9.2}{
  \begin{scope}[xshift=\shift cm,x=1.5cm,y=1.1cm]
   \draw[->] (-1.1,0)--(1.12,0) node[right]{$\delta_1$};
   \draw[->] (0,-1.5)--(0,1.52) node[above]{$\delta_2$};
   \draw[densely dotted,thick] (-.737210,-1.105816) rectangle (.737210,1.105816);
   \draw[domain=0:360,samples=100,smooth,variable=\t]
   plot ({.816114*cos(\t)},
    {1.224171*(\corr*cos(\t)+sqrt(1-\corr*\corr)*sin(\t))});
   \node at (0,-1.9) {$\rho=\corr$};
  \end{scope}
 }
\end{tikzpicture}
\caption{Giữ phương sai thành phần, tâm và thang trục giống nhau:
 nét chấm Bonferroni không đổi; nét liền elip thay đổi theo tương quan.
 Các trục là độ lệch khỏi trung bình mẫu.}
\label{fig:rho_regions}
\end{figure}

\hieu{Hình chữ nhật trả lời ``mỗi trung bình có thể ở khoảng nào?''.
Elip còn trả lời ``hai trung bình có thể ở những vị trí ấy \emph{cùng
nhau} không?''. Với tương quan dương, cặp lệch cùng chiều được đánh giá
khác cặp lệch ngược chiều, dù từng độ lệch riêng có cùng độ lớn.}

Tuy nhiên, không suy rằng ``Bonferroni hoàn toàn không liên quan tương
quan''. Cần phân biệt ba điều:
\begin{enumerate}
 \item \textbf{Công thức biên của hình chữ nhật thành phần} không dùng
 hiệp phương sai chéo. Hình chữ nhật từ các khoảng thành phần $T^2$
 cũng có tính chất này nếu giữ đường chéo, $n,p,\alpha$ cố định.
 \item \textbf{Xác suất phủ đồng thời thực tế} có thể phụ thuộc tương
 quan vì các thống kê một chiều phụ thuộc nhau. Bonferroni chỉ bảo đảm
 một cận dưới, không khẳng định xác suất phủ bằng nhau trong mọi mô hình.
 \item \textbf{Khoảng cho tổ hợp} vẫn dùng $\mathbf a'S\mathbf a$.
 Với hiệu, $\Var$ mẫu là $s_{11}+s_{22}-2s_{12}$; tương quan dương
 làm khoảng hiệu hẹp hơn khi giữ các phương sai thành phần cố định.
\end{enumerate}

\paragraph{Nếu thêm một hiệu thì miền không còn đơn thuần chữ nhật.}
Khi danh sách có $\mu_1,\mu_2,\mu_1-\mu_2$, phải dùng $m=3$ cho
toàn bộ danh sách. Hai khoảng thành phần tạo hình chữ nhật rộng hơn
trường hợp $m=2$. Khoảng hiệu thêm điều kiện
\[
 |(\mu_1-\mu_2)-(\bar x_1-\bar x_2)|
 \le t_{n-1;1-\alpha/6}\sqrt{(s_{11}+s_{22}-2s_{12})/n}.
\]
Đó là một dải chéo; giao của dải với hình chữ nhật thường là đa giác.
Trong ví dụ hiện tại, điều kiện dải là $|\delta_1-\delta_2|\le1.0465$.
Không vẽ riêng hình chữ nhật rồi nói nó thể hiện đầy đủ cả ba khoảng.

\kiemtra{Giữ $s_{11},s_{22}$ và $n$ nhưng đổi dấu $s_{12}$:
 hình chữ nhật thành phần Bonferroni có đổi không? Không.
 Elip có đổi không? Có, hướng nghiêng đổi.
 Khoảng cho $\mu_1-\mu_2$ có đổi không? Có, vì có hạng $-2s_{12}$.}


\subsection{Khi nào Bonferroni hẹp hơn?}
Cho cùng một tổ hợp, cùng dữ liệu, cùng mức $\alpha$, tỉ số độ dài là
\begin{equation}
 \frac{\text{độ dài Bonferroni}}{\text{độ dài }T^2}
 =\frac{t_{n-1;1-\alpha/(2m)}}{\sqrt{c_T^2}}.
 \label{eq:widthratio}
\end{equation}
Chỉ khi tỉ số nhỏ hơn 1 mới nói Bonferroni hẹp hơn. Khi $m$ tăng, hệ số
Bonferroni tăng; do đó kết luận không đúng cho mọi danh sách hữu hạn.
Ví dụ với $n=100,p=2,\alpha=0.05$, hệ số $T^2$ khoảng 2.4983,
còn với $m=10000$, hệ số Bonferroni khoảng 4.8291.

Điều quan trọng không chỉ là độ rộng: phải chọn bảo đảm phù hợp câu hỏi.
Nếu chỉ có vài phát biểu xác định trước, Bonferroni dễ áp dụng.
Nếu cần mọi hướng hoặc chọn hướng sau khi xem dữ liệu, $T^2$ phù hợp hơn.
Không tính nhiều khoảng rồi chỉ giữ những khoảng ``đẹp'' mà coi $m$ là
số khoảng cuối cùng đã chọn.

\kiemtra{Các khoảng Bonferroni có cần độc lập không? Không.
Nếu có 4 trung bình và 1 hiệu cần báo cáo thì $m=5$, không phải 4.}
